\section{Fishbowl Assessment Procedure}

Here's how we will conduct the \href{https://www.learningforjustice.org/classroom-resources/teaching-strategies/community-inquiry/fishbowl}{fishbowl} assessment. Each group will be given approximately 30 minutes to discuss about how you group apply the principle of innovation process and product development while developing the selected product. 

The purpose of this activity is not only to assess your individual performance but also to ensure that the group collaborates effectively to create a viable product. During the discussion, other students will observe your group. As the name (fishbowl)  of the assessment imply, the setup will be similar to a fishbowl, where your group will be in the center and the others will watch from around the sides.

Please note that this assessment is not about testing your English language skills but rather about how effectively you can communicate your ideas, how well you understand and contribute to the overall success of the group's product.

As a community, everyone naturally wants to see their classmates succeed. Therefore, during the fishbowl assessment, your peers from other groups, who will be observing from outside the central discussion area, will engage actively. They will not only ask constructive questions to challenge your ideas and thinking but will also offer valuable suggestions on how to improve your group’s project. This collaborative approach helps ensure that all students have the opportunity to refine their work and achieve better results.
Following each fishbowl session, we will divide into smaller groups to review the performance. Individual feedback will be given based on each group member's contributions during the fishbowl assessment. This discussion will help clarify the strengths and areas for improvement, ensuring that everyone receives personalized and constructive feedback to foster their development.


\subsection{Assessment Criteria}

Your individual performance during the fishbowl will be judged using the following criteria:

\begin{itemize}
    \item \textbf{Body Language \& Facial Expressions} – Demonstrates attentiveness and engagement through non-verbal communication.
    \item \textbf{Confidence \& Engagement} – Actively participates and speaks with clarity and assurance.
    \item \textbf{Defense of Ideas} – Provides logical and well-supported justifications for the group’s decisions.
    \item \textbf{Connection to Group Objectives} – Aligns individual points with the broader goals of the group project.
    \item \textbf{Critical Thinking} – Shows depth in analysis, responds thoughtfully to questions, and considers alternative perspectives.
\end{itemize}

\subsection{Preparation Tips}
\begin{itemize}
    \item review your team's current needs, specifications, and concept choices before the session;
    \item prepare short explanations for important technical decisions;
    \item bring notes or supporting visuals if the lecturer allows them; and
    \item listen carefully to questions so that your answers remain relevant and precise.
\end{itemize}

\subsection{Reflection Prompt}

After the fishbowl, note which questions exposed weak reasoning or missing evidence in your project. Use that feedback to strengthen the final presentation and report.
